% Shared layout for the four manuals. Compile with XeLaTeX (see docs/tools/build_manuals.py).
\usepackage[a4paper,margin=2.3cm,headheight=14pt]{geometry}
\usepackage{xcolor}
\usepackage{graphicx}
\usepackage{booktabs,tabularx,array,longtable}
\usepackage{enumitem}
\usepackage{fancyhdr}
\usepackage{titlesec}
\usepackage{amsmath,amssymb}
\usepackage{listings}
\usepackage[most]{tcolorbox}
\usepackage[hidelinks]{hyperref}

\definecolor{ink}{HTML}{1F2937}
\definecolor{accent}{HTML}{0F5C8C}
\definecolor{soft}{HTML}{EEF3F7}
\definecolor{warnbg}{HTML}{FFF6E5}
\definecolor{warnline}{HTML}{C47F00}
\definecolor{codebg}{HTML}{F7F7F5}
\definecolor{codegray}{HTML}{8A8F98}
\definecolor{codekw}{HTML}{0F5C8C}
\definecolor{codestr}{HTML}{8A4B08}
\definecolor{codecom}{HTML}{4E7A3E}

\color{ink}
\setlist{itemsep=2pt,topsep=4pt}
\setlength{\parindent}{0pt}
\setlength{\parskip}{5pt}

\titleformat{\section}{\Large\bfseries\color{accent}}{\thesection}{0.8em}{}
\titleformat{\subsection}{\large\bfseries}{\thesubsection}{0.7em}{}
\titleformat{\subsubsection}{\normalsize\bfseries}{\thesubsubsection}{0.6em}{}
\titleformat{\part}[display]{\centering\Huge\bfseries\color{accent}}{\partname~\thepart}{0.5em}{}

\pagestyle{fancy}
\fancyhf{}
\renewcommand{\headrulewidth}{0.3pt}
\fancyfoot[C]{\small\thepage}

\lstset{
  language=Python,
  basicstyle=\ttfamily\footnotesize,
  backgroundcolor=\color{codebg},
  keywordstyle=\color{codekw}\bfseries,
  stringstyle=\color{codestr},
  commentstyle=\color{codecom},
  numbers=left,
  numberstyle=\tiny\color{codegray},
  numbersep=6pt,
  frame=single,
  rulecolor=\color{codegray!40},
  breaklines=true,
  breakatwhitespace=true,
  postbreak=\mbox{\textcolor{codegray}{$\hookrightarrow$}\space},
  columns=fullflexible,
  keepspaces=true,
  showstringspaces=false,
  tabsize=4,
  upquote=true,
  xleftmargin=1.6em,
  framexleftmargin=1.4em,
  aboveskip=8pt,
  belowskip=8pt,
  captionpos=t,
}
\lstdefinestyle{shell}{language=bash,numbers=none,xleftmargin=0.6em,framexleftmargin=0.4em,
  keywordstyle=\color{ink},commentstyle=\color{codegray}}
\lstdefinestyle{plain}{language={},numbers=none,xleftmargin=0.6em,framexleftmargin=0.4em}
\lstnewenvironment{shell}{\lstset{style=shell}}{}
\lstnewenvironment{plaincode}{\lstset{style=plain}}{}

% Result numbers (\RlinEst etc.) come from the latest pipeline run.
\InputIfFileExists{generated/results.tex}{}{\typeout{Run docs/tools/build_manuals.py first}}

% \snippet{path:name} prints the current source of a function or class.
\InputIfFileExists{generated/snippets.tex}{}{\typeout{Run docs/tools/build_manuals.py first}}
\newcommand{\snippet}[1]{%
  \ifcsname snip@#1\endcsname\csname snip@#1\endcsname
  \else\fbox{\ttfamily missing snippet: \detokenize{#1}}\fi}

\newtcolorbox{note}[1][Note]{colback=soft,colframe=accent,boxrule=0.5pt,arc=1pt,
  left=6pt,right=6pt,top=4pt,bottom=4pt,fonttitle=\bfseries\small,title=#1,breakable}
\newtcolorbox{warn}[1][Watch out]{colback=warnbg,colframe=warnline,boxrule=0.5pt,arc=1pt,
  left=6pt,right=6pt,top=4pt,bottom=4pt,fonttitle=\bfseries\small,title=#1,breakable}

\newcommand{\file}[1]{\texttt{\detokenize{#1}}}
\newcommand{\cmd}[1]{\texttt{\detokenize{#1}}}
\newcolumntype{L}{>{\raggedright\arraybackslash}X}
